% =====================================================================
%  Chapter 26 — Decision Criteria
% =====================================================================

\chapter{Candidate-by-Candidate Decision Criteria}
\label{ch:criteria}

This chapter converts the survey into something a funding body or a
laboratory could act on. For each candidate class: what result would
confirm it, what would exclude it, what is the realistic route, how long
would it take, and what would justify stopping.

\section{The criteria}Five questions are asked of every candidate, and they are
asked of every candidate in the same form so that the answers can be
compared.

\begin{enumerate}
\item \textbf{Confirmation.} What observation, at what significance, in
which experiment would make this candidate the accepted answer?
\item \textbf{Exclusion.} What result would eliminate it, and is that
result achievable in the foreseeable future?
\item \textbf{Route.} What is the cheapest path to either answer, given
that the obvious one has failed?
\item \textbf{Scale.} What is the timescale, and what does the answer cost?
\item \textbf{Kill condition.} What observation would justify abandoning the
line of work, and who decides?
\end{enumerate}

The fifth question is the one the field is worst at, and its absence is why
some lines of work persist for decades without producing knowledge. A
candidate that cannot be excluded by any foreseeable observation should not
be funded as a candidate; it should be funded as a study of a model class,
which is a different activity with different criteria.

\section{The comparison}The table below applies the five questions to the
candidate classes of Part~II. The timescale column is the time to a
\emph{decisive} result either way, not the time to a result at all.

\begin{table}[!htbp]
\centering
\caption{Decision criteria by candidate class. The scale column is the time to
a decisive result either way, not the time to a result at all. The fuller
reasoning for each row is in the chapter named in the text above.}
\label{tab:criteria}
\setlength{\tabcolsep}{4pt}
\footnotesize
\begin{tabular}{@{}>{\RaggedRight\arraybackslash}p{0.150\textwidth}p{0.200\textwidth}p{0.175\textwidth}p{0.090\textwidth}p{0.180\textwidth}@{}}
\toprule
\textbf{Candidate} & \textbf{Confirmation} & \textbf{Route} & \textbf{Scale} & \textbf{Kill condition} \\
\midrule
WIMP & $5\sigma$, two experiments & light targets, electron recoils & 10--20~yr & a null over the whole thermal region \\
Axion & photon line, two haloscopes & haloscopes, helioscopes & 10--20~yr & the canonical band covered at full sensitivity \\
Ultralight field & lensing phase; soliton & CMB lensing, dwarf dynamics & 5--10~yr & small-scale power seen where none is allowed \\
Primordial black hole & microlensing or merger & wide-field surveys, GW data & 5--10~yr & survey depth of a few per cent \\
SIMP & cores in dwarf galaxies & dwarf dynamics & 10~yr & limits reaching the SIMP region \\
Freeze-in & almost nothing directly & cosmology only & --- & --- \\
Modified gravity & --- & cluster lensing, CMB & 5~yr & a completion that fits the lensing offset \\
Sterile neutrino & an X-ray or UV line & X-ray spectroscopy & 10~yr & line searches covering the region \\
Superheavy & a beam-dump event & beam dumps & 5--10~yr & sensitivity over the claimed range \\
\bottomrule
\end{tabular}
\end{table}

\tabref{tab:criteria} is the operational core of this book. Three features
of it are worth drawing out.

First, the confirmation column is not ``a signal''. It is a signal, at a
significance, in more than one experiment, with a specific structure. That
requirement is what Chapter~\ref{ch:event248} is about, and it should be
applied to the 248~keV candidate exactly as it is applied to everything
else.

Second, the freeze-in row has no timescale and no kill condition, and that
is the honest entry for a candidate whose defining property is that it
cannot be detected. A candidate that cannot be excluded should be funded as
a model study rather than as a search, and it should be funded at a level
proportional to its explanatory value rather than to its popularity.

Third, the modified-gravity row is short on timescale because the exclusion
is already effectively complete at the cluster scale. That is a cheap
exclusion of an expensive hypothesis, and it is the best return in the
table.

\section{How this table should be used}Three cautions.

The kill conditions are deliberately stated as observations rather than as
budget decisions. A line of work should be stopped when the physics says
so, and the physics in this field has repeatedly surprised people who
assumed their candidate was safe.

The timescales assume the current funding rate and the current
instrumentation pipeline, and they are estimates rather than
commitments. The 2025--2026 experience with the 248~keV candidate is the
relevant caution: a single event can change the field's priorities in
weeks, and every timescale in this table should be read as conditional on
no such event occurring.

And the table omits the candidates that are not yet candidates --- the ones
that have not been proposed, or that the current instruments cannot address.
That omission is the main risk in any such table, and it is the reason
Chapter~\ref{ch:roadmap} reserves part of the programme for
exploratory work.

\begin{keytakeaways}
\item Every candidate class is assessed against the same five questions,
including the one the field avoids: what would justify stopping.
\item The table exposes candidates that cannot be excluded and should
therefore be funded as model studies rather than as searches.
\item The cheapest exclusion in the table --- modified gravity --- is already
effectively complete.
\item The table's blind spot is the candidates not yet proposed, which is
why the roadmap reserves exploratory effort.
\end{keytakeaways}
